\documentclass[10pt,CJKmath]{ctexart}
\usepackage{MyNote}




% 定义重点高亮和公式框样式
\usepackage{framed,longtable,booktabs,array,calc}

\usepackage{etoolbox}

\AtBeginEnvironment{longtable}{\footnotesize}
\AtEndEnvironment{longtable}{\normalsize}


%参考tcolorbox P21
\NewTotalTCBox{\commandbox}{ s v }
{verbatim,colupper=white,colback=black!50!white,colframe=black,sharp corners,boxrule=0pt,left=0mm,right=0mm,top=0mm,bottom=0mm}
{\IfBooleanT{#1}{\textcolor{red}{\ttfamily\bfseries > }}%
	\lstinline[language=c++,basicstyle=\ttfamily\small,keywordstyle=\color{blue!35!white}\bfseries]^#2^}

\usepackage{pifont}
%使用带圈的数字0-10，过多则会出错
\newcommand{\makenum}[1]{{\ding{\numexpr171+#1}}}%

\newenvironment{myenum}{
\begin{dingautolist}{182}\setlength{\itemsep}{0pt}\setlength{\parskip}{0pt}%
}
{\end{dingautolist}}
	
\begin{document}

\pagestyle{fancynotes}


\part{Attribute Set}

\section{\texttt{PreAttributeChange}}

\textbf{核心结论：}
$\begin{cases}
	\text{CurrentValue}\Rightarrow\text{PreAttributeChange} \\
	\text{BaseValue}\Rightarrow\text{PreAttributeBaseChange}
\end{cases}$

\begin{myenum}
	\item \textbf{两者触发时机}

		\begin{itemize}
			\item \ci{PreAttributeChange}在属性的 \textbf{当前值（CurrentValue）} 即将变化时触发。  
			常见于带持续时间/无限时长 GE 的修饰器引起的数值变化（聚合后更新 Current）。
			
			\item \ci{PreAttributeBaseChange}
			在属性的 \textbf{基础值（BaseValue）} 即将变化时触发。  
			常见于 Instant GE 或直接改 Base 的路径（如 \texttt{SetNumericAttributeBase}）。
		\end{itemize}

\item \textbf{为什么会有两个钩子？}

在 GAS 里通常可理解为：\colorbox{yellow!5}{$
\text{CurrentValue} = \mathbf{f}(BaseValue, Modifiers)
$}
因此：
\begin{itemize}
	\item 改 \textbf{BaseValue} 前，用 \texttt{PreAttributeBaseChange}；
	\item 改 \textbf{CurrentValue} 前，用 \texttt{PreAttributeChange}。
\end{itemize}

\item \textbf{实战建议（很重要）}

\begin{itemize}
	\item 如果你要做夹逼（clamp），例如生命值限制在 \([0, MaxHealth]\)，
	最稳妥做法是 \textbf{两个函数都处理}，避免不同修改路径漏掉。
	\item 对 Instant GE 的“最终结果”校正，通常还会配合下面这个收尾钩子：
	\begin{itemize}
		\item \ci{PostGameplayEffectExecute}
	\end{itemize}
\end{itemize}
\end{myenum}

在 GAS 中，带 Period 的 GE 每次周期触发会按“执行(Execute)”语义处理，这条路径本质上接近 Instant GE，通常会动到 BaseValue。 而\textbf{一旦 BaseValue 变了，聚合器会重新计算 CurrentValue}，于是你通常会看到：
\begin{itemize}
	\item \ci{PreAttributeBaseChange(...)}：Base 将要改；
	\item \ci{PreAttributeChange(...)}：Current 随后重算将要改。
\end{itemize}
所以“两个都调用”是正常现象，不是异常。

可以把流程理解成：\colorbox{yellow}{$
\text{Periodic Tick} \Rightarrow \Delta \text{BaseValue} \Rightarrow \text{Recompute CurrentValue}$}。

因此：\ci{PreAttributeBaseChange} 先触发，\ci{PreAttributeChange} 后触发。

\paragraph{为什么你会觉得“按道理只改 Current”？}
因为“持续时间/无限时长 GE”这类描述通常让人联想到“临时 Buff/Debuff（Modifier 挂在 Active GE 上）”，那条路径主要表现为 Current 层变化。
但一旦你加了 Period，每个周期会走一次执行逻辑，常常落到 Base 改动路径上，所以就把两边都带起来了。

实战建议（Aura 项目里很实用）
\begin{enumerate}
	\item \ci{PreAttributeBaseChange}：放“基础值合法化”逻辑（如防 NaN、下限保护）。
	\item \ci{PreAttributeChange}：放“显示值 / 运行时值的 clamp”（如 Health 限制到 [0, MaxHealth]）。
	\item \ci{PostGameplayEffectExecute}：处理“执行后副作用”（死亡判定、受击、溢出转护盾等）。
	\item 避免在两个 \ci{Pre} 里写同一套重逻辑，否则容易重复处理。
\end{enumerate}

\section{PostGameplayEffectExecute}
\subsection{基础信息}
该函数使用参数\ci{struct FGameplayEffectModCallbackData}，定义如下：
\begin{code}
struct FGameplayEffectModCallbackData
{
	`$\dots$`
	const struct FGameplayEffectSpec&		EffectSpec;		// The spec that the mod came from
	struct FGameplayModifierEvaluatedData&	EvaluatedData;	// The 'flat'/computed data to be applied to the target
	
	class UAbilitySystemComponent &Target;		// Target we intend to apply to
};
\end{code}
可以从此结构中获取许多信息：哪个GE触发这个函数？效果应用在谁身上\sn{Target引用}？及最重要的\ci{EvaluatedData}，影响哪些数值。
\subsubsection{详解ATTRIBUTE\_ACCESSORS\_BASIC宏}
当使用宏：\ci{ATTRIBUTE_ACCESSORS_BASIC(ThisClass, Health)}时，会展开成如下代码（以下做了适当简化，避免干扰）：
\begin{code}[caption={简化的宏展开}]

static FGameplayAttribute GetHealthAttribute()
{
	static FProperty* Prop = FindFieldChecked<FProperty>(ThisClass::StaticClass(),
	FName("Health"));
	return Prop;
}

inline float GetHealth() const { return Health.GetCurrentValue(); }

inline void SetHealth(float NewVal)
{
	UAbilitySystemComponent* AbilityComp = GetOwningAbilitySystemComponent();
	AbilityComp->SetNumericAttributeBase(GetHealthAttribute(), NewVal);
}

inline void InitHealth(float NewVal)
{
	Health.SetBaseValue(NewVal);
	Health.SetCurrentValue(NewVal);
}
\end{code}
这个宏其实内联生成4个函数，分别对应属性及值的\ci{Init/Get/Set}函数。尤其要注意\ci{SetHealth}，它设置的是\textbf{Base值}。别忘记上节所述，Base值变动会引起Current值的变动\sn{Current值相当于$f(BaseValue,Modifiers)$}。
\begin{code}[caption={Base值变动，触发所有当前激活的GE重新设置Base值}]
void UAbilitySystemComponent::SetNumericAttributeBase(const FGameplayAttribute &Attribute, float NewFloatValue)
{
	// Go through our active gameplay effects container so that aggregation/mods are handled properly.
	ActiveGameplayEffects.SetAttributeBaseValue(Attribute, NewFloatValue);
}	
\end{code}



\subsection{1. 调用时机}

\ci{PostGameplayEffectExecute}
是在 \texttt{GameplayEffect} \textbf{已经执行完并且属性值已经被修改之后} 调用的回调。

可简单理解为：
	GameplayEffect 改值完成 $\Rightarrow$ \ci{PostGameplayEffectExecute}

它通常对应：
\begin{itemize}
	\item Instant GE
	\item Periodic GE 的每次周期结算
	\item ExecutionCalculation 最终写入属性之后
\end{itemize}
\begin{marker}
	 Duration GE 不触发 \ci{PostGameplayEffectExecute}
\end{marker}
\subsection{2. 适合做什么}

这个函数最适合做“\textbf{最终结算}”：

\begin{itemize}
	\item 对属性做 Clamp
	\item 处理伤害、治疗、蓝量结算
	\item 死亡判定
	\item 根据 \verb|Data| 读取来源、目标、标签等上下文
\end{itemize}

例如：
\begin{itemize}
	\item \verb|Health| 改变后限制到 \verb|[0, MaxHealth]|
	\item \verb|Mana| 改变后限制到 \verb|[0, MaxMana]|
\end{itemize}

\subsection{3. 与其他回调的区别}

\subsubsection{(1) 与 \texttt{PreAttributeChange} 的区别}

\begin{itemize}
	\item \verb|PreAttributeChange|：属性\textbf{将要}变化时调用
	\item \ci{PostGameplayEffectExecute}：属性\textbf{已经}变化后调用
\end{itemize}

所以：

\begin{itemize}
	\item \verb|PreAttributeChange| 更像“前置处理”
	\item \ci{PostGameplayEffectExecute} 更像“后置结算”
\end{itemize}

\subsubsection{(2) 与 \texttt{OnRep\_XXX} 的区别}

\begin{itemize}
	\item \ci{PostGameplayEffectExecute}：偏服务端结算逻辑
	\item \verb|OnRep_Health| / \verb|OnRep_Mana|：偏客户端同步与 UI 刷新
\end{itemize}

即：

\begin{quote}
	\ci{PostGameplayEffectExecute} 负责“规则结算”，  
	\verb|OnRep| 负责“复制后响应”。
\end{quote}

\subsection{4. 常见写法}

最常见的写法是根据本次被修改的属性分支处理：

\begin{code}
if (Data.EvaluatedData.Attribute == GetHealthAttribute())
{
	// 处理 Health
}
	
if (Data.EvaluatedData.Attribute == GetManaAttribute())
{
	// 处理 Mana
}
\end{code}

含义是：

\begin{quote}
	检查这次 GameplayEffect 实际改动的是哪个属性，然后分别做对应结算。
\end{quote}

\subsection{5. 注意事项}

\subsubsection{(1) 不是所有属性变化都会走这里}

\textbf{不要把它理解成“任意属性变化通知”。}

它主要是在 \verb|GameplayEffect| 执行并写入属性时触发。  
如果属性是通过别的方式修改的，不一定会调用这里。

\subsubsection{(2) 更适合做服务端权威逻辑}

常见放在这里的逻辑有：

\begin{itemize}
	\item Clamp
	\item 扣血/回血
	\item 死亡判定
	\item 资源修正
\end{itemize}

不建议把大量 UI 或表现层逻辑直接堆在这里。

\subsubsection{(3) 小心重复结算}

如果你在这里又施加新的 GE，或者再次修改相关属性，要注意：

\begin{itemize}
	\item 是否会重复触发
	\item 是否会形成递归
	\item 是否会导致一次伤害被算两遍
\end{itemize}

\subsubsection{(4) 推荐用于最终兜底}

对于生命值、蓝量这类有上下限的属性，  
\ci{PostGameplayEffectExecute} 很适合做最终 Clamp，保证结果合法。

\subsection{6. 一句话总结}

\begin{quote}
	\textbf{\ci{PostGameplayEffectExecute} 就是 GameplayEffect 改完属性后的最终结算点。}
\end{quote}

它最适合做：

\begin{itemize}
	\item Clamp
	\item 伤害/治疗收尾
	\item 死亡与资源判定
\end{itemize}

使用时要注意：

\begin{itemize}
	\item 它不是所有属性变化的统一入口
	\item 它更偏向服务端数值逻辑
	\item 客户端 UI 更新通常应结合 \verb|OnRep|
\end{itemize}

\subsection{对 PostGameplayEffectExecute 的准确理解}

可以近似理解为：

\begin{quote}
	\ci{PostGameplayEffectExecute} 是 \textbf{GameplayEffect 执行型结算之后} 的回调。
\end{quote}

它的触发关键并不是“属性改了 \verb|BaseValue|”本身，而是：

\begin{center}
	\textbf{该次属性修改是否来自 GameplayEffect 的 Execute 路径}
\end{center}

在 GAS 中，典型会触发它的情况有：

\begin{itemize}
	\item Instant GameplayEffect
	\item Periodic GameplayEffect 的每次周期执行
	\item ExecutionCalculation 的最终属性写入
\end{itemize}

这些情况通常都会修改属性的 \verb|BaseValue|，因此可以得到一个常见现象：

\begin{quote}
	\textbf{凡是触发 PostGameplayEffectExecute 的，大多数时候都伴随着 \ci{BaseValue} 的修改。}
\end{quote}

但反过来不能机械理解为：

\begin{quote}
	“只要 \verb|BaseValue| 变了，就一定会触发 \ci{PostGameplayEffectExecute}。”
\end{quote}

更不能把它当成“所有属性变化的统一通知入口”。

特别是：

\begin{itemize}
	\item \verb|Duration| / \verb|Infinite| 类型的 GameplayEffect 通常只是通过 Modifier 影响 \verb|CurrentValue|
	\item 这类变化一般不会调用 \ci{PostGameplayEffectExecute}
\end{itemize}

因此，一句话总结为：

\begin{quote}
	\textbf{PostGameplayEffectExecute 本质上是 Execute 型 GameplayEffect 的结算后回调，而不是单纯的 BaseValue 变化回调。}
\end{quote}

\subsection{什么是 GameplayEffect 的 Execute 路径}

所谓 \textbf{Execute 路径}，可以理解为：

\begin{quote}
	\textbf{GameplayEffect 不是仅仅挂上一个持续 Modifier，而是实际执行了一次数值结算，并将结果真正写入属性。}
\end{quote}

其核心过程可概括为：

\begin{center}
	计算修改量 $\rightarrow$ 执行属性写入 $\rightarrow$ 触发结算后回调
\end{center}

在 GAS 中，\ci{PostGameplayEffectExecute} 主要就是在这种“\textbf{执行型结算完成之后}”被调用的。

\subsubsection{1. Execute 路径的直观理解}

可以把 Execute 路径理解成“\textbf{结算落账}”：

\begin{itemize}
	\item 火球命中，立刻扣除 30 点生命
	\item 治疗命中，立刻恢复 20 点生命
	\item 周期伤害每跳一次，真正扣除一次生命
\end{itemize}

这类变化的特点是：

\begin{itemize}
	\item 它是一次真实的数值事件
	\item 结算结果会沉淀到属性中
	\item 通常对应属性 \verb|BaseValue| 语义上的变化
	\item 通常会触发 \ci{PostGameplayEffectExecute}
\end{itemize}

\subsubsection{2. 与 Modifier 路径的区别}

与 Execute 路径相对的是“\textbf{持续 Modifier 生效路径}”，它更像“\textbf{状态挂账}”：

\begin{itemize}
	\item 角色身上挂了一个 Buff，于是当前生命临时 +50
	\item 身上挂了一个 Debuff，于是当前护甲临时 -10
	\item 光环存在期间，移速持续乘以 1.2
\end{itemize}

这类变化更像：

\begin{quote}
	效果存在时，属性的最终生效值被持续影响；效果消失时，影响随之消失。
\end{quote}

因此它通常：

\begin{itemize}
	\item 影响的是 \verb|CurrentValue|
	\item 不属于一次“执行结算写入”
	\item 一般不会调用 \ci{PostGameplayEffectExecute}
\end{itemize}

\subsection{Execute 路径与 GE Duration Policy 的对应关系}

\begin{table}[h]
	\centering
	\begin{longtable}{|p{0.1\textwidth}|p{0.1\textwidth}|p{0.2\textwidth}|p{0.35\textwidth}|p{0.25\textwidth}|}
		\hline
		\textbf{GE 类型} & \textbf{是否属于 Execute 路径} & \textbf{常见影响目标} & \textbf{是否常触发 \texttt{PostGameplayEffectExecute}} & \textbf{典型用途} \\
		\hline
		Instant & 是 & 通常是 \verb|BaseValue| 语义 & 是 & 立即伤害、立即治疗、立即回蓝、一次性资源变化 \\
		\hline
		Duration & 否（通常） & 通常影响 \verb|CurrentValue| & 否（通常） & 持续 Buff / Debuff，例如持续加攻击、持续减防御 \\
		\hline
		Infinite & 否（通常） & 通常影响 \verb|CurrentValue| & 否（通常） & 装备加成、常驻光环、持续被动效果 \\
		\hline
		Periodic & 是（每个周期执行时） & 每次周期结算通常是 \verb|BaseValue| 语义 & 是（每次周期执行时） & 中毒每跳扣血、持续回蓝每跳加蓝、灼烧每秒伤害 \\
		\hline
	\end{longtable}
	\caption{GameplayEffect 的 Execute 路径与持续 Modifier 路径对比}
\end{table}

\subsection{为什么说它“通常对应 BaseValue”，但又不能简单等同}

很多时候可以观察到：\colorbox{yellow}{
	$触发  \texttt{PostGameplayEffectExecute} \Rightarrow 常常伴随 \text{BaseValue} 变化$}

因此容易形成一种简化理解：

\begin{quote}
	“只有修改 \verb|BaseValue| 的才会触发 \ci{PostGameplayEffectExecute}。”
\end{quote}

这个理解\textbf{接近正确}，但更准确的说法应当是：

\begin{quote}
	\textbf{\ci{PostGameplayEffectExecute} 的触发关键不是“BaseValue 变了”，而是“GameplayEffect 走了 Execute 路径”。}
\end{quote}

也就是说：

\begin{center}
	\textbf{Execute} $\Rightarrow$ \ci{PostGameplayEffectExecute}
\end{center}

而不是机械地写成：

\begin{center}
	\verb|BaseValue| \text{变化} $\Rightarrow$ \ci{PostGameplayEffectExecute}
\end{center}

\subsection{为什么 Duration 型 GE 不适合直接处理伤害结算}

如果一个 \commandbox{Duration GE} 只是给 \verb|Health| 挂了一个持续的负 Modifier，例如：

\[
CurrentHealth = BaseHealth + (-50)
\]

那么会出现两个问题：

\begin{enumerate}
	\item 它通常不会触发 \ci{PostGameplayEffectExecute}
	\item 效果消失时，生命值可能“弹回去”
\end{enumerate}

这会让“死亡判定”变得很不稳定，例如：

\begin{itemize}
	\item Debuff 生效时，目标当前生命低于 0
	\item Debuff 消失后，又恢复到正数
\end{itemize}

因此在战斗系统里，更常见也更稳妥的做法是：

\begin{itemize}
	\item \textbf{伤害}：使用 Instant 或 Periodic Execute 进行真实扣血
	\item \textbf{持续伤害}：使用 Periodic，每次 tick 都执行一次真实扣血
	\item \textbf{Buff / Debuff}：使用 Duration / Infinite 影响攻击、防御、速度、护盾等
\end{itemize}

\subsection{与 \texttt{PostGameplayEffectExecute} 的关系}

\ci{PostGameplayEffectExecute} 可以理解为：

\begin{quote}
	\textbf{Execute 型 GameplayEffect 完成属性结算后，AttributeSet 拿到的一次后置处理机会。}
\end{quote}

因此它特别适合做：

\begin{itemize}
	\item \verb|Health| Clamp
	\item \verb|Mana| Clamp
	\item 死亡判定
	\item 护盾扣减后的血量修正
	\item 伤害、治疗、资源变化的最终收尾
\end{itemize}

但不应把它理解成：

\begin{itemize}
	\item 所有属性变化的统一通知
	\item 所有 \verb|CurrentValue| 变化都会触发的回调
	\item 客户端 UI 更新的唯一入口
\end{itemize}

\subsection{结合 AuraAttributeSet 的理解}

在当前的 `UAuraAttributeSet` 中，\ci{PostGameplayEffectExecute} 被用于：

\begin{itemize}
	\item 当 \verb|Health| 被 Execute 型效果修改后，限制在 \verb|[0, MaxHealth]|
	\item 当 \verb|Mana| 被 Execute 型效果修改后，限制在 \verb|[0, MaxMana]|
\end{itemize}

这说明当前代码的设计思路是：

\begin{quote}
	\textbf{把 Health 和 Mana 当作需要“真实结算落地”的属性，并在结算后统一做兜底修正。}
\end{quote}

如果后续伤害系统继续沿着这个思路扩展，那么推荐做法是：

\begin{itemize}
	\item 用 Instant / Periodic Execute 处理扣血与回血
	\item 在 \ci{PostGameplayEffectExecute} 中处理 Clamp、死亡、受击收尾
	\item 不要用 Duration 直接给 \verb|Health| 挂持续负 Modifier 来表示真实伤害
\end{itemize}

\subsection{一句话总结}

\begin{quote}
	\textbf{Execute 路径 = GameplayEffect 真正执行了一次数值结算并把结果写入属性；而 \ci{PostGameplayEffectExecute} 就是这次结算之后的后置回调。}
\end{quote}
\part{Gameplay Tags}
目标:在ASC中监听GE应用到目标Actor上，然后根据GE的AssetTag显示UI。监听事件使用委托：\ci{OnGameplayEffectAppliedDelegateToSelf}，使用\ci{Asset Tag}标识GE。在UI显示时需要使用一些额外的信息，例如GE的名称，效果之类，这样的信息可以统一保存在数据表中。这个数据表行结构用\cppsign{}定义，这就是本节内容。
\section{UI Widget数据表}
行结构要从\ci{FTableRowBase}派生：
\begin{code}[caption={数据表}]
USTRUCT(BlueprintType)
struct FUIWidgetRow : public FTableRowBase
{
	GENERATED_BODY()
	
	UPROPERTY(EditAnywhere, BlueprintReadOnly)
	FGameplayTag MessageTag;
	
	UPROPERTY(EditAnywhere, BlueprintReadOnly)
	FText Message;
	
	UPROPERTY(EditAnywhere, BlueprintReadOnly)
	TSubclassOf<UAuraUserWidget> MessageWidgetClass;
	
	UPROPERTY(EditAnywhere, BlueprintReadOnly)
	UTexture2D* Image;
};
\end{code}
\ci{OnGameplayEffectAppliedDelegateToSelf}是Server only函数，注定在客户端不会触发，对于UI要显示，需要使用别的方法。

\section{Attribute Based Modifier 的基本用法与作用}

本节介绍了 Gameplay Effect 中另一种重要的数值计算方式：\texttt{Attribute Based}。此前课程中的 Modifier Magnitude 一直使用 \texttt{Scalable Float}，而 \texttt{Attribute Based} 的核心价值在于，它可以让一个属性的变化量直接基于其他属性的当前值来计算，因此非常适合实现“派生属性”或“由其他属性驱动的属性变化”。

为了演示这一点，课程先新建了一个测试用的 \texttt{Effect Actor}，并为其添加一个可见的 Box Collision，在角色进入碰撞体时调用 \texttt{ApplyEffectToTarget}。随后又创建了一个测试用的 Gameplay Effect，命名为 \texttt{GE\_Test\_Attribute\_Based}，并将其设置为该测试 Actor 的 Instant Gameplay Effect。这样，Aura 角色一旦进入该区域，就会立即受到这个效果的影响。

接着，课程重点讲解了如何配置 \texttt{Attribute Based} Modifier。示例中给 Health 添加一个 Modifier，操作类型保持为 \texttt{Add}，但将 Magnitude Calculation Type 从 \texttt{Scalable Float} 改为 \texttt{Attribute Based}。在其展开配置中，\texttt{Backing Attribute} 用来指定“以哪个属性为基础”进行计算，\texttt{Attribute Source} 则决定这个属性是从 Source 还是 Target 身上获取。由于测试场景中是将效果施加给 Aura，所以这里选择从 \texttt{Target} 捕获属性，并以 Aura 的 \texttt{Vigor} 作为基础属性。这样，当角色进入测试区域时，Health 就会增加当前 Vigor 的数值；示例中 Aura 的 Health 为 50、Vigor 为 9，因此进入后 Health 变为 59。课程还简要提到 \texttt{Snapshot} 用于决定何时捕获属性值，但这一机制会留到后面更合适的场景再深入讨论。

之后，课程进一步演示了多个 Modifier 同时作用于同一属性的情况。除了第一个基于 Vigor 的 Health 加成外，又添加了第二个同样作用于 Health 的 Modifier，这次改为基于 \texttt{Strength}。由于两个 Modifier 的操作类型都是 \texttt{Add}，最终效果就是把 Vigor 和 Strength 的值都加到 Health 上。示例中角色初始 Health 为 50，Vigor 为 9，Strength 为 10，因此进入测试区域后 Health 变为 69，结果符合直觉。

本节的结论是：\texttt{Attribute Based} 提供了一种非常灵活的方式，使 Gameplay Effect 的数值不再是固定常数，而是可以由角色自身或来源对象的属性动态驱动。当前示例展示的是最简单、最直观的加法情况，而更复杂的行为还涉及多个 Modifier 叠加时的执行顺序，以及 \texttt{Coefficient}、\texttt{Pre-Multiply Additive}、\texttt{Post-Multiply Additive} 等附加参数。下一节将继续深入分析不同 Modifier Operation 下的运算顺序与具体规则。

\section{SynchronizeProperties 的作用}

在 UMG / \texttt{UUserWidget} 体系中，\texttt{SynchronizeProperties()} 的核心职责可以概括为：

\begin{quote}
	把对象上保存的属性值，同步到实际的 Slate / Widget 显示层。
\end{quote}

它更偏向于“属性同步”而不是“业务逻辑初始化”。

很多时候可以把它理解为：

\begin{itemize}
	\item C++ 成员变量里已经有了一些值；
	\item 现在需要把这些值真正应用到界面控件上；
	\item 这一步就是 \texttt{SynchronizeProperties()} 要做的事情。
\end{itemize}

例如：

\begin{itemize}
	\item 把一个 \texttt{FText TitleText} 同步到 \texttt{TextBlock->SetText(...)}；
	\item 把一个 \texttt{float BoxWidth} 同步到 \texttt{SizeBox->SetWidthOverride(...)}；
	\item 把一个颜色变量同步到 \texttt{Border} 或 \texttt{Image} 上。
\end{itemize}

\section{它和 NativePreConstruct / NativeConstruct 的区别}

这三个函数都和 Widget 初始化有关，但关注点不同。

\subsection{NativePreConstruct}

\texttt{NativePreConstruct()} 更偏向：

\begin{itemize}
	\item 在设计器预览和运行时都可能调用；
	\item 用来提前调整 UI 的外观；
	\item 适合做“展示层的预配置”。
\end{itemize}

例如：

\begin{itemize}
	\item 设置默认文本；
	\item 设置尺寸；
	\item 设置可见性；
	\item 根据可编辑属性刷新外观。
\end{itemize}

\subsection{NativeConstruct}

\texttt{NativeConstruct()} 更偏向：

\begin{itemize}
	\item 运行时逻辑初始化；
	\item 绑定按钮事件；
	\item 获取游戏对象；
	\item 注册委托；
	\item 创建动态子控件。
\end{itemize}

它通常不是专门为“属性同步”设计的，而是为“控件开始参与游戏逻辑”准备的。

\subsection{SynchronizeProperties}

\texttt{SynchronizeProperties()} 则更聚焦于：

\begin{itemize}
	\item 将当前对象的属性值应用到实际控件；
	\item 保证界面显示和对象状态一致；
	\item 更像“刷新显示层”。
\end{itemize}

因此可以把它和前两个函数粗略区分成：

\begin{itemize}
	\item \texttt{NativePreConstruct()}：预构造阶段调整外观；
	\item \texttt{NativeConstruct()}：运行时接入逻辑；
	\item \texttt{SynchronizeProperties()}：把属性值同步到界面。
\end{itemize}

\section{什么时候会调用}

从使用者角度，不必死记底层所有细节，但需要知道一个重要事实：

\begin{quote}
	\texttt{SynchronizeProperties()} 不是“业务初始化入口”，而是“属性刷新入口”。
\end{quote}

它通常会在引擎认为“需要把对象属性重新应用到界面”时被调用。  
这类场景一般包括：

\begin{itemize}
	\item Widget 构建或重建时；
	\item 设计器预览需要刷新时；
	\item 某些属性变化后需要重新同步时。
\end{itemize}

所以它可能不是只调用一次，而可能被多次调用。

这也是为什么在这个函数里通常应该做：

\begin{itemize}
	\item 幂等的 UI 同步；
	\item 轻量的显示刷新；
	\item 不依赖复杂运行时上下文的操作。
\end{itemize}

\section{适合放什么代码}

\subsection{适合放的内容}

适合放进 \texttt{SynchronizeProperties()} 的，一般是“把已有属性值写到控件上”的代码，例如：

\begin{itemize}
	\item 文本同步；
	\item 颜色同步；
	\item 图片同步；
	\item 尺寸同步；
	\item 可见性同步；
	\item 某些简单状态的表现层刷新。
\end{itemize}

例如：

\begin{code}
	if (TitleTextBlock)
	{
		TitleTextBlock->SetText(TitleText);
	}
	
	if (SizeBox_Root)
	{
		SizeBox_Root->SetWidthOverride(BoxWidth);
		SizeBox_Root->SetHeightOverride(BoxHeight);
	}
\end{code}

这里的思想是：

\begin{quote}
	成员变量是“数据”，Widget 子控件是“显示”，\\
	\texttt{SynchronizeProperties()} 负责把前者写到后者。
\end{quote}

\subsection{不适合放的内容}

不建议放进 \texttt{SynchronizeProperties()} 的内容包括：

\begin{itemize}
	\item 获取 \texttt{PlayerController} / \texttt{GameMode} / \texttt{PlayerState}；
	\item 绑定按钮点击事件；
	\item 注册委托；
	\item 创建大量动态子控件；
	\item 依赖运行时世界状态的复杂逻辑；
	\item 带副作用的业务操作。
\end{itemize}

原因是：

\begin{enumerate}
	\item 这个函数可能被多次调用；
	\item 它的职责本来就是“同步属性”，不是“做业务初始化”；
	\item 如果在这里绑定事件或重复创建对象，很容易导致重复执行。
\end{enumerate}

\section{一个典型使用模式}

一个常见做法是：

\begin{enumerate}
	\item 在 \ci{UPROPERTY(EditAnywhere)} 里暴露一些可编辑参数；
	\item 在 \ci{SynchronizeProperties()} 中把这些参数应用到子控件；
	\item 这样在设计器和运行时都更容易保持显示一致。
\end{enumerate}

例如一个 Widget 有如下属性：

\begin{code}
	UPROPERTY(EditAnywhere, BlueprintReadWrite, Category="UI")
	float BoxWidth = 805.f;
	
	UPROPERTY(EditAnywhere, BlueprintReadWrite, Category="UI")
	float BoxHeight = 906.f;
\end{code}

那么就可以在 \texttt{SynchronizeProperties()} 中写：

\begin{code}
	void UAttributeMenu::SynchronizeProperties()
	{
		Super::SynchronizeProperties();
		
		if (SizeBox_Root)
		{
			SizeBox_Root->SetWidthOverride(BoxWidth);
			SizeBox_Root->SetHeightOverride(BoxHeight);
		}
	}
\end{code}

这样它的语义非常明确：

\begin{quote}
	每当需要同步属性时，就把 \texttt{BoxWidth} 和 \texttt{BoxHeight} 应用到 \texttt{SizeBox\_Root}。
\end{quote}

\section{和你的 AttributeMenu 场景如何对应}

从你当前的 \texttt{`AttributeMenu.h`} 来看，你已经有：

\begin{itemize}
	\item \texttt{SizeBox\_Root}
	\item \texttt{BoxWidth}
	\item \texttt{BoxHeight}
	\item \texttt{UpdateFrameSize()}
\end{itemize}

这类“根据可编辑属性同步 UI 尺寸”的行为，其实就很适合使用 \texttt{SynchronizeProperties()}。

可以把思路理解为：

\begin{itemize}
	\item \texttt{BoxWidth} / \texttt{BoxHeight} 是数据；
	\item \texttt{SizeBox\_Root} 是显示对象；
	\item \texttt{SynchronizeProperties()} 就是把这两个数据写入显示对象。
\end{itemize}

因此如果 \texttt{UpdateFrameSize()} 的职责仅仅是：

\begin{itemize}
	\item 调用 \texttt{SetWidthOverride()}；
	\item 调用 \texttt{SetHeightOverride()}；
\end{itemize}

那么它就非常适合在 \texttt{SynchronizeProperties()} 中调用。

例如逻辑上可以写成：

\begin{code}
	void UAttributeMenu::SynchronizeProperties()
	{
		Super::SynchronizeProperties();
		UpdateFrameSize();
	}
\end{code}

\section{为什么它在做“编辑器友好型 UI”时很有用}

当一个 Widget 需要在编辑器里就能看到配置效果时，\texttt{SynchronizeProperties()} 特别有价值。

因为这种场景里，你经常会有这样的需求：

\begin{itemize}
	\item 改一个 \texttt{UPROPERTY(EditAnywhere)}；
	\item 立即让对应控件的显示变化反映出来；
	\item 保持设计器预览和运行时显示尽量一致。
\end{itemize}

这正是“属性同步函数”最擅长做的事情。

所以对于：

\begin{itemize}
	\item 面板尺寸；
	\item 默认标题；
	\item 颜色主题；
	\item 图标样式；
	\item 文本样式；
\end{itemize}

这样的可编辑展示参数，\texttt{SynchronizeProperties()} 通常是很合适的切入点。

\section{使用时的注意事项}

\subsection{要调用父类实现}

通常应当先调用：

\begin{code}
	Super::SynchronizeProperties();
\end{code}

这样可以保证父类已有的同步逻辑正常执行。

\subsection{尽量保持幂等}

因为它可能被多次调用，所以最好保证：

\begin{itemize}
	\item 调一次和调多次结果一致；
	\item 不产生额外副作用；
	\item 不重复绑定事件；
	\item 不重复创建对象。
\end{itemize}

\subsection{注意空指针判断}

由于子控件可能在某些阶段尚未准备好，通常要先判断：

\begin{code}
	if (SizeBox_Root)
	{
		...
	}
\end{code}

\subsection{不要把它当成业务入口}

如果你需要：

\begin{itemize}
	\item 和角色属性系统通信；
	\item 绑定按钮事件；
	\item 初始化运行时数据源；
	\item 构建滚动列表内容；
\end{itemize}

这些通常更适合放到 \texttt{NativeConstruct()} 或专门的初始化函数中，而不是 \texttt{SynchronizeProperties()}。

\section{一个简明判断标准}

可以用下面这个原则区分：

\begin{quote}
	如果你的代码是在说“把这个属性显示到界面上”，就优先考虑 \texttt{SynchronizeProperties()}；\\
	如果你的代码是在说“让这个界面开始参与游戏逻辑”，就优先考虑 \texttt{NativeConstruct()}。
\end{quote}

\section{一句话总结}

\texttt{SynchronizeProperties()} 的本质是：

\begin{quote}
	\textbf{把 Widget 对象上的配置属性，同步到实际界面表现。}
\end{quote}

它适合做轻量、可重复、偏显示层的刷新；  
而不适合承担复杂的运行时初始化和业务逻辑。
\end{document}